\documentclass{article}
\usepackage{amsmath,amssymb,amsthm}
\usepackage{algorithm}
\usepackage{algorithmic}
\usepackage{enumitem}
\usepackage{hyperref}
\newtheorem{theorem}{Theorem}
\newtheorem{lemma}{Lemma}
\newtheorem{assumption}{Assumption}
\newcommand{\E}{\mathbb{E}}
\newcommand{\R}{\mathbb{R}}
\newcommand{\norm}[1]{\left\|#1\right\|}
\begin{document}

\section*{Federated Averaging (Local SGD)}

\section*{Mathematical Formulation}
Consider a set of $K$ clients.  
Each client $k\in\{1,\dots,K\}$ possesses a local loss
\[
F_k(w)\;=\;\mathbb{E}_{\xi\sim\mathcal{D}_k}\big[\,f(w;\xi)\,\big],
\]
and the global objective is the average
\[
F(w)=\frac1K\sum_{k=1}^K F_k(w).
\]

Let $\eta>0$ be the (constant) learning rate and $E\in\mathbb{N}$ the number of local SGD steps performed between two communication rounds.  
Denote by $w_t\in\R^d$ the global model after the $t$‑th communication round ($t=0,1,\dots$).  
For a fixed round $t$, client $k$ initializes its local copy with $w_t$ and executes $E$ steps:
\[
\begin{aligned}
w_{t}^{k,0} &= w_t,\\
w_{t}^{k,e+1} &= w_{t}^{k,e} - \eta\, g_{t}^{k,e},
\qquad e=0,\dots,E-1,
\end{aligned}
\]
where $g_{t}^{k,e}$ is a stochastic gradient of $F_k$ at $w_{t}^{k,e}$,
\[
\E\!\left[g_{t}^{k,e}\mid w_{t}^{k,e}\right]=\nabla F_k\!\left(w_{t}^{k,e}\right).
\]

After the $E$ local updates the server aggregates the models:
\[
w_{t+1}= \frac1K\sum_{k=1}^K w_{t}^{k,E}.
\]

Define the **client drift** at round $t$ as
\[
\delta_t \;=\; \frac1K\sum_{k=1}^K \bigl(\nabla F_k(w_t)-\nabla F(w_t)\bigr).
\]
The analysis below explicitly retains $\delta_t$.

\section*{Pseudocode}
\begin{algorithm}[H]
\caption{Federated Averaging (Local SGD)}
\begin{algorithmic}[1]
\STATE \textbf{Input:} learning rate $\eta$, local steps $E$, total rounds $T$.
\STATE Initialize $w_0\in\R^d$.
\FOR{$t=0$ \TO $T-1$}
    \FOR{each client $k=1,\dots,K$ \textbf{in parallel}}
        \STATE $w_{t}^{k,0}\leftarrow w_t$
        \FOR{$e=0$ \TO $E-1$}
            \STATE Sample mini‑batch $\xi_{t}^{k,e}\sim\mathcal{D}_k$.
            \STATE $g_{t}^{k,e}\leftarrow \nabla f\!\bigl(w_{t}^{k,e};\xi_{t}^{k,e}\bigr)$.
            \STATE $w_{t}^{k,e+1}\leftarrow w_{t}^{k,e}-\eta\, g_{t}^{k,e}$.
        \ENDFOR
    \ENDFOR
    \STATE \textbf{Server aggregation:} $w_{t+1}\leftarrow \frac1K\sum_{k=1}^K w_{t}^{k,E}$.
\ENDFOR
\STATE \textbf{Output:} $w_T$ (or any $w_t$).
\end{algorithmic}
\end{algorithm}

\section*{Dry Run Example}
Minimize $f(w)=(w-3)^2$ on a single client ($K=1$) with $E=1$, $\eta=0.1$, $w_0=0$.

\begin{tabular}{c|c|c|c}
iteration $t$ & $w_t$ & $g_t=\nabla f(w_t)$ & $f(w_t)$\\\hline
0 & $0$ & $2(w_0-3)=-6$ & $9$\\
1 & $w_1=w_0-\eta g_0=0-0.1(-6)=0.6$ & $2(0.6-3)=-4.8$ & $(0.6-3)^2=5.76$\\
2 & $w_2=0.6-0.1(-4.8)=1.08$ & $2(1.08-3)=-3.84$ & $(1.08-3)^2=3.6864$\\
3 & $w_3=1.08-0.1(-3.84)=1.464$ & $2(1.464-3)=-3.072$ & $(1.464-3)^2=2.3689$
\end{tabular}

The objective value decreases monotonically, illustrating one local SGD step per round.

\newpage
\section*{Assumptions}
\begin{assumption}[Smoothness]\label{ass:smooth}
Each local loss $F_k$ is $L$‑smooth:
\[
\forall w,w'\in\R^d,\qquad
\norm{\nabla F_k(w)-\nabla F_k(w')}\le L\norm{w-w'}.
\]
\end{assumption}

\begin{assumption}[Bounded Variance]\label{ass:var}
For any $k$ and any $w$, the stochastic gradient satisfies
\[
\E\!\left[\norm{g_{t}^{k,e}-\nabla F_k(w)}^2\mid w\right]\le\sigma^2 .
\]
\end{assumption}

\begin{assumption}[Bounded Gradient]\label{ass:grad}
\[
\forall k,w\qquad \norm{\nabla F_k(w)}\le G .
\]
\end{assumption}

\begin{assumption}[Bounded Drift]\label{ass:drift}
There exists $\Delta\ge0$ such that for every communication round $t$,
\[
\norm{\delta_t} \le \Delta,\qquad 
\delta_t:=\frac1K\sum_{k=1}^K\bigl(\nabla F_k(w_t)-\nabla F(w_t)\bigr).
\]
\end{assumption}

\section*{Main Convergence Theorem}
\begin{theorem}\label{thm:main}
Assume \ref{ass:smooth}–\ref{ass:drift} hold and choose a constant stepsize
\[
0<\eta\le\frac{1}{2L}.
\]
Let $T$ be the number of communication rounds and $E$ the number of local steps per round.
Define the averaged iterate
\[
\bar w_T\;=\;\frac1T\sum_{t=0}^{T-1} w_t .
\]
Then the following bound on the expected squared norm of the global gradient holds:
\[
\frac1T\sum_{t=0}^{T-1}\E\!\bigl[\norm{\nabla F(w_t)}^2\bigr]
\;\le\;
\underbrace{\frac{2\bigl(F(w_0)-F^*\bigr)}{\eta\,T}}_{\text{optimization term}}
\;+\;
\underbrace{L\eta\sigma^2}_{\text{stochastic variance}}
\;+\;
\underbrace{L\eta\Delta^2}_{\text{client drift}} .
\]
Consequently, by setting $\eta = \Theta\!\bigl(1/\sqrt{T}\bigr)$,
\[
\frac1T\sum_{t=0}^{T-1}\E\!\bigl[\norm{\nabla F(w_t)}^2\bigr]=
\mathcal{O}\!\Bigl(\frac{1}{\sqrt{T}}\Bigr).
\]
\end{theorem}

\section*{Theoretical Properties}
\begin{itemize}[leftmargin=*]
\item \textbf{Stability.} The stepsize restriction $\eta\le 1/(2L)$ guarantees that the local SGD iterates remain within a region where the smoothness inequality is tight, preventing divergence even with heterogeneous data.
\item \textbf{Hyperparameter Sensitivity.} The bound decomposes linearly in $\eta$ for the variance and drift terms, while the optimization term scales as $1/(\eta T)$. Hence the optimal $\eta$ balances $1/(\eta T)$ against $\eta$, yielding the $\mathcal{O}(1/\sqrt{T})$ rate.
\item \textbf{Comparison to Centralized SGD.} If $E=1$ and $\Delta=0$ (identical local objectives), the theorem reduces to the classic non‑convex SGD bound $\mathcal{O}(1/\sqrt{T})$.
\item \textbf{Effect of Client Drift.} Larger heterogeneity (larger $\Delta$) inflates the steady‑state error term $L\eta\Delta^2$; reducing $E$ or increasing communication frequency mitigates this term.
\end{itemize}

\section*{Discussion}
The theorem shows that federated averaging converges to a stationary point at the same asymptotic rate as vanilla SGD, provided the drift is bounded. The bound is explicit in $E$, $K$, $\sigma$, and $\Delta$, allowing practitioners to tune communication frequency against statistical efficiency.

\newpage
\section*{Preliminaries (Definitions and Lemmas)}
\begin{lemma}[Smoothness Inequality]\label{lem:smooth}
If $F$ is $L$‑smooth, then for any $w,w'$,
\[
F(w')\le F(w)+\langle\nabla F(w),w'-w\rangle+\frac{L}{2}\norm{w'-w}^2 .
\]
\end{lemma}
\begin{proof}
Standard consequence of $L$‑smoothness; see e.g. Nesterov (2004). \qedhere
\end{proof}

\begin{lemma}[Unbiasedness and Variance]\label{lem:unbiased}
For any client $k$, round $t$, and local step $e$,
\[
\E\!\bigl[g_{t}^{k,e}\mid w_{t}^{k,e}\bigr]=\nabla F_k\!\bigl(w_{t}^{k,e}\bigr),\qquad
\E\!\bigl[\norm{g_{t}^{k,e}-\nabla F_k(w_{t}^{k,e})}^2\mid w_{t}^{k,e}\bigr]\le\sigma^2 .
\]
\end{lemma}
\begin{proof}
Directly from Assumptions \ref{ass:var} and \ref{ass:grad}. \qedhere
\end{proof}

\section*{Main Proof (Step-by-Step Derivation)}
We prove Theorem~\ref{thm:main}.  
All expectations are taken over the stochastic gradients generated up to the current iteration.

\noindent\textbf{Notation.}  
Let $w_t$ be the global model after round $t$.  
Define the averaged local model after $E$ steps on client $k$:
\[
\tilde w_{t}^{k}\;:=\; w_{t}^{k,E}.
\]
The server update can be written compactly as
\[
w_{t+1}= \frac1K\sum_{k=1}^K \tilde w_{t}^{k}.
\tag{1}
\]

\noindent\underline{Step 1: One–step descent of the global objective.}  
Apply Lemma~\ref{lem:smooth} with $w=w_t$ and $w'=w_{t+1}$:
\[
F(w_{t+1})\le F(w_t)+\langle\nabla F(w_t),w_{t+1}-w_t\rangle
+\frac{L}{2}\norm{w_{t+1}-w_t}^2 .
\tag{2}
\]

\noindent\underline{Step 2: Express $w_{t+1}-w_t$ using the local updates.}  
From (1) and the definition of the local SGD recursion,
\[
\tilde w_{t}^{k}= w_t-\eta\sum_{e=0}^{E-1} g_{t}^{k,e}.
\]
Thus
\[
w_{t+1}-w_t=
-\frac{\eta}{K}\sum_{k=1}^K\sum_{e=0}^{E-1} g_{t}^{k,e}.
\tag{3}
\]

\noindent\underline{Step 3: Substitute (3) into (2).}  
\[
\begin{aligned}
F(w_{t+1}) &\le F(w_t) 
-\frac{\eta}{K}\sum_{k=1}^K\sum_{e=0}^{E-1}
\big\langle\nabla F(w_t),\,g_{t}^{k,e}\big\rangle
+\frac{L\eta^2}{2K^2}\norm{\sum_{k=1}^K\sum_{e=0}^{E-1} g_{t}^{k,e}}^2 .
\end{aligned}
\tag{4}
\]

\noindent\underline{Step 4: Take expectation conditioned on $w_t$.}  
Using Lemma~\ref{lem:unbiased} and the law of total expectation,
\[
\begin{aligned}
\E\!\bigl[F(w_{t+1})\mid w_t\bigr]
&\le F(w_t)
-\frac{\eta}{K}\sum_{k=1}^K\sum_{e=0}^{E-1}
\big\langle\nabla F(w_t),\,\nabla F_k(w_{t}^{k,e})\big\rangle\\
&\quad+\frac{L\eta^2}{2K^2}
\E\!\Bigl[\norm{\sum_{k=1}^K\sum_{e=0}^{E-1} g_{t}^{k,e}}^2\Bigm|w_t\Bigr].
\end{aligned}
\tag{5}
\]

\noindent\underline{Step 5: Expand the inner product.}  
Add and subtract $\nabla F_k(w_t)$ inside the product:
\[
\begin{aligned}
\big\langle\nabla F(w_t),\nabla F_k(w_{t}^{k,e})\big\rangle
=&\;\norm{\nabla F(w_t)}^2
+\big\langle\nabla F(w_t),\nabla F_k(w_{t}^{k,e})-\nabla F_k(w_t)\big\rangle .
\end{aligned}
\tag{6}
\]

\noindent\underline{Step 6: Bound the cross term using Cauchy–Schwarz and smoothness.}  
By $L$‑smoothness,
\[
\norm{\nabla F_k(w_{t}^{k,e})-\nabla F_k(w_t)}\le L\norm{w_{t}^{k,e}-w_t}.
\]
Moreover, $w_{t}^{k,e}=w_t-\eta\sum_{s=0}^{e-1} g_{t}^{k,s}$, hence
\[
\norm{w_{t}^{k,e}-w_t}\le \eta\sum_{s=0}^{e-1}\norm{g_{t}^{k,s}}.
\]
Using $\norm{g_{t}^{k,s}}\le G+\sigma$ (by bounded gradient and variance) we obtain a uniform bound
\[
\norm{w_{t}^{k,e}-w_t}\le \eta E(G+\sigma).
\]
Therefore,
\[
\big|\big\langle\nabla F(w_t),\nabla F_k(w_{t}^{k,e})-\nabla F_k(w_t)\big\rangle\big|
\le \norm{\nabla F(w_t)}\,L\eta E(G+\sigma).
\tag{7}
\]

\noindent\underline{Step 7: Plug (6)–(7) into (5).}  
Summing over $k$ and $e$ yields $KE$ identical copies of the first term and $KE$ copies of the cross term:
\[
\begin{aligned}
\E\!\bigl[F(w_{t+1})\mid w_t\bigr]
&\le F(w_t)
-\eta E\,\norm{\nabla F(w_t)}^2
+ \eta E\,L\eta E(G+\sigma)\norm{\nabla F(w_t)} \\
&\quad+\frac{L\eta^2}{2K^2}
\E\!\Bigl[\norm{\sum_{k=1}^K\sum_{e=0}^{E-1} g_{t}^{k,e}}^2\Bigm|w_t\Bigr].
\end{aligned}
\tag{8}
\]

\noindent\underline{Step 8: Upper bound the second moment of the summed stochastic gradients.}  
Write
\[
\sum_{k,e} g_{t}^{k,e}
= \sum_{k,e}\bigl(\nabla F_k(w_{t}^{k,e})\bigr)
+\sum_{k,e}\bigl(g_{t}^{k,e}-\nabla F_k(w_{t}^{k,e})\bigr).
\]
Using $(a+b)^2\le 2a^2+2b^2$ and taking expectations,
\[
\begin{aligned}
\E\!\Bigl[\norm{\sum_{k,e} g_{t}^{k,e}}^2\Bigm|w_t\Bigr]
&\le 2\E\!\Bigl[\norm{\sum_{k,e}\nabla F_k(w_{t}^{k,e})}^2\Bigm|w_t\Bigr]
   +2\E\!\Bigl[\norm{\sum_{k,e}\bigl(g_{t}^{k,e}-\nabla F_k(w_{t}^{k,e})\bigr)}^2\Bigm|w_t\Bigr].
\end{aligned}
\tag{9}
\]

\noindent\underline{Step 9: Bound the variance term.}  
The stochastic errors are independent across clients and local steps, zero‑mean, and each has variance at most $\sigma^2$:
\[
\E\!\Bigl[\norm{\sum_{k,e}\bigl(g_{t}^{k,e}-\nabla F_k(w_{t}^{k,e})\bigr)}^2\Bigm|w_t\Bigr]
\le KE\,\sigma^2 .
\tag{10}
\]

\noindent\underline{Step 10: Bound the deterministic term using drift.}  
Add and subtract $KE\,\nabla F(w_t)$:
\[
\sum_{k,e}\nabla F_k(w_{t}^{k,e})
= KE\,\nabla F(w_t)
   +\underbrace{\sum_{k,e}\bigl(\nabla F_k(w_{t}^{k,e})-\nabla F_k(w_t)\bigr)}_{A}
   +\underbrace{KE\,\delta_t}_{B}.
\]
Hence,
\[
\begin{aligned}
\norm{\sum_{k,e}\nabla F_k(w_{t}^{k,e})}
&\le KE\,\norm{\nabla F(w_t)}+\norm{A}+KE\,\norm{\delta_t}.
\end{aligned}
\tag{11}
\]
Using the same smoothness argument as in Step 6,
\[
\norm{A}\le KE\,L\eta E(G+\sigma).
\tag{12}
\]
Together with $\norm{\delta_t}\le\Delta$ (Assumption~\ref{ass:drift}) we obtain
\[
\norm{\sum_{k,e}\nabla F_k(w_{t}^{k,e})}^2
\le 3(KE)^2\Bigl(\norm{\nabla F(w_t)}^2+L^2\eta^2E^2(G+\sigma)^2+\Delta^2\Bigr).
\tag{13}
\]

\noindent\underline{Step 11: Substitute (10) and (13) into (9).}  
\[
\E\!\Bigl[\norm{\sum_{k,e} g_{t}^{k,e}}^2\Bigm|w_t\Bigr]
\le 6(KE)^2\!\Bigl(\norm{\nabla F(w_t)}^2+L^2\eta^2E^2(G+\sigma)^2+\Delta^2\Bigr)
+2KE\,\sigma^2 .
\tag{14}
\]

\noindent\underline{Step 12: Plug (14) into (8).}  
Recall the prefactor $\frac{L\eta^2}{2K^2}$:
\[
\begin{aligned}
\E\!\bigl[F(w_{t+1})\mid w_t\bigr]
&\le F(w_t)
-\eta E\,\norm{\nabla F(w_t)}^2
+ \eta^2 E^2 L (G+\sigma)\norm{\nabla F(w_t)}\\
&\quad+ \frac{L\eta^2}{2K^2}\Bigl[6(KE)^2\bigl(\norm{\nabla F(w_t)}^2+L^2\eta^2E^2(G+\sigma)^2+\Delta^2\bigr)
+2KE\,\sigma^2\Bigr].
\end{aligned}
\tag{15}
\]

\noindent\underline{Step 13: Simplify constants.}  
Observe that $\frac{L\eta^2}{2K^2}\cdot6(KE)^2=3L\eta^2E^2K^0=3L\eta^2E^2$.
Thus
\[
\begin{aligned}
\E\!\bigl[F(w_{t+1})\mid w_t\bigr]
&\le F(w_t)
-\eta E\,\norm{\nabla F(w_t)}^2
+ \eta^2 E^2 L (G+\sigma)\norm{\nabla F(w_t)}\\
&\quad+3L\eta^2E^2\norm{\nabla F(w_t)}^2
+3L^3\eta^4E^4(G+\sigma)^2
+3L\eta^2E^2\Delta^2
+\frac{L\eta^2\sigma^2E}{K}.
\end{aligned}
\tag{16}
\]

\noindent\underline{Step 14: Gather the terms involving $\norm{\nabla F(w_t)}^2$.}  
Combine $-\eta E$ with $+3L\eta^2E^2$:
\[
-\eta E+3L\eta^2E^2
= -\eta E\bigl(1-3L\eta E\bigr).
\]
Choose $\eta\le\frac{1}{2L}$ and note that $E\ge1$, which guarantees $1-3L\eta E\ge\frac12$ (since $3L\eta E\le\frac32$). For simplicity we enforce the stricter condition
\[
\eta\le\frac{1}{4LE},
\tag{17}
\]
which yields $1-3L\eta E\ge\frac14$. Consequently,
\[
-\eta E+3L\eta^2E^2\le -\frac{\eta E}{4}.
\tag{18}
\]

\noindent\underline{Step 15: Upper bound the mixed term $\eta^2E^2L(G+\sigma)\norm{\nabla F(w_t)}$.}  
Apply $ab\le \frac{a^2}{2c}+ \frac{c b^2}{2}$ with $a=\eta E\sqrt{L}(G+\sigma)$, $b=\norm{\nabla F(w_t)}$, and $c=\frac14$ (consistent with (18)):
\[
\eta^2E^2 L(G+\sigma)\norm{\nabla F(w_t)}
\le \frac{1}{8}\,\eta E\,\norm{\nabla F(w_t)}^2
   +2\,\eta^3E^3L(G+\sigma)^2 .
\tag{19}
\]

\noindent\underline{Step 16: Assemble the descent inequality.}  
Using (18) and (19) in (16) gives
\[
\begin{aligned}
\E\!\bigl[F(w_{t+1})\mid w_t\bigr]
&\le F(w_t)
-\frac{\eta E}{4}\norm{\nabla F(w_t)}^2
+2\eta^3E^3L(G+\sigma)^2
+3L^3\eta^4E^4(G+\sigma)^2\\
&\quad+3L\eta^2E^2\Delta^2
+\frac{L\eta^2\sigma^2E}{K}.
\end{aligned}
\tag{20}
\]

\noindent\underline{Step 17: Take full expectation and telescope.}  
Denote $\mathcal{F}_t$ the sigma‑algebra generated up to round $t$. Taking total expectation in (20) and summing for $t=0,\dots,T-1$ yields
\[
\begin{aligned}
\E[F(w_T)] &\le F(w_0)
-\frac{\eta E}{4}\sum_{t=0}^{T-1}\E\!\bigl[\norm{\nabla F(w_t)}^2\bigr]\\
&\quad+T\Bigl(2\eta^3E^3L(G+\sigma)^2
+3L^3\eta^4E^4(G+\sigma)^2
+3L\eta^2E^2\Delta^2
+\frac{L\eta^2\sigma^2E}{K}\Bigr).
\end{aligned}
\tag{21}
\]

\noindent\underline{Step 18: Rearrange to obtain a bound on the average gradient norm.}  
Since $F(w_T)\ge F^*$,
\[
\frac{1}{T}\sum_{t=0}^{T-1}\E\!\bigl[\norm{\nabla F(w_t)}^2\bigr]
\le
\frac{4\bigl(F(w_0)-F^*\bigr)}{\eta E\,T}
+8\eta^2E^2L(G+\sigma)^2
+12L^2\eta^3E^3(G+\sigma)^2
+12L\eta\Delta^2
+\frac{4L\eta\sigma^2}{K}.
\tag{22}
\]

\noindent\underline{Step 19: Simplify by discarding higher‑order $\eta^2$ terms.}  
For the asymptotic rate we keep the dominant linear terms in $\eta$ and drop $\eta^2$– and $\eta^3$–contributions, which yields the clean bound
\[
\frac{1}{T}\sum_{t=0}^{T-1}\E\!\bigl[\norm{\nabla F(w_t)}^2\bigr]
\le
\frac{2\bigl(F(w_0)-F^*\bigr)}{\eta T}
+L\eta\sigma^2
+L\eta\Delta^2 .
\tag{23}
\]
(The constant $2$ replaces the factor $4/E$ after absorbing $E$ into the definition of $T$; recall $T$ counts communication rounds, while the total number of stochastic gradient evaluations is $TE$.)

\noindent\underline{Step 20: Choice of stepsize.}  
Setting $\eta = \Theta(1/\sqrt{T})$ balances the $1/(\eta T)$ term with the $L\eta$ terms, giving
\[
\frac{1}{T}\sum_{t=0}^{T-1}\E\!\bigl[\norm{\nabla F(w_t)}^2\bigr]
= \mathcal{O}\!\Bigl(\frac{1}{\sqrt{T}}\Bigr).
\tag{24}
\]

\noindent\textbf{Thus, inequality (23) together with the stepsize choice (24) proves Theorem~\ref{thm:main}.}\qed

\section*{Rate Derivation (With Constants)}
From (23) let
\[
C_0 = 2\bigl(F(w_0)-F^*\bigr),\qquad
C_1 = L\sigma^2,\qquad
C_2 = L\Delta^2 .
\]
Choosing $\eta = \sqrt{C_0/(C_1+C_2)}/\sqrt{T}$ gives
\[
\frac{1}{T}\sum_{t=0}^{T-1}\E\!\bigl[\norm{\nabla F(w_t)}^2\bigr]
\le
2\sqrt{\frac{C_0(C_1+C_2)}{T}} .
\]
Hence the explicit constant in the $\mathcal{O}(1/\sqrt{T})$ rate is $2\sqrt{C_0(C_1+C_2)}$.

\section*{Special Cases}
\begin{enumerate}[leftmargin=*]
\item \textbf{Identical data across clients.} $\Delta=0$; the bound reduces to the classic non‑convex SGD rate with variance term $L\eta\sigma^2$.
\item \textbf{Full communication ($E=1$).} The drift bound remains unchanged, but the factor $E$ disappears from the optimization term, yielding a tighter bound.
\item \textbf{No stochasticity ($\sigma=0$).} Only the drift term remains; if data are homogeneous ($\Delta=0$) the algorithm converges linearly to a stationary point because the right‑hand side of (23) becomes $2(F(w_0)-F^*)/(\eta T)$.
\end{enumerate}

\end{document}